\documentclass[aspectratio=169]{beamer}

% Tema base
\usetheme{metropolis}
\usecolortheme{default}

% Palette colori 
\definecolor{thesisPrimary}{HTML}{BB2E29}
\definecolor{thesisSecondary}{HTML}{53565A}

\setbeamercolor{palette primary}{bg=thesisPrimary, fg=white}
\setbeamercolor{palette secondary}{bg=thesisSecondary, fg=white}
\setbeamercolor{block title}{bg=thesisPrimary, fg=white}
\setbeamercolor{block body}{bg=thesisPrimary!8, fg=black}
\setbeamercolor{section in toc}{fg=thesisPrimary}
\setbeamercolor{subsection in toc}{fg=thesisPrimary}
\setbeamercolor{item projected}{bg=thesisPrimary, fg=white}

% Pacchetti
\usepackage[T1]{fontenc}
\usepackage[utf8]{inputenc}
\usepackage[italian]{babel}
\usepackage{lmodern}
\usepackage{listings}
\usepackage{xcolor}
\usepackage{graphicx}
\usepackage{tikz}
\usepackage{booktabs}

% Metadati
\title{Progettazione di un framework protocol-agnostic per il coordinamento di sciami nel contesto di Project Emerge}
\author{Alessandro Gardini}
\institute{Corso di Laurea in Ingegneria e Scienze Informatiche \\ Relatore: Mirko Viroli \\ Correlatori: Gianluca Aguzzi, Nicolas Farabegoli}
\date{A.A. 2025/2026}

\setbeamertemplate{footline}[frame number]
\setbeamertemplate{navigation symbols}{}
\AtBeginSection[]{}

\begin{document}
%
\begin{frame}
    \titlepage
\end{frame}
%
\begin{frame}{Indice}
    \tableofcontents
\end{frame}
%
\section{Introduzione}
%
\section{Background}
%
\begin{frame}{Background: Aggregate Computing}
    \begin{columns}
        \column{0.45\textwidth}
        \centering
        Paradigma per programmare sciami di dispositivi eterogenei come un'unica entità collettiva, tramite \textbf{campi computazionali} -- senza controllo centralizzato
        
        \column{0.03\textwidth}

        \column{0.45\textwidth}
        \begin{block}{Fondamenta e strumenti}
            Formalizzato dal \textit{field calculus}; su questa base: \textbf{MacroSwarm} (libreria di comportamenti di sciame) e \textbf{ScaFi} (framework di esecuzione in Scala)
        \end{block}
    \end{columns}

    \vspace{2.1em}

    \begin{block}{Esecuzione}
            Round asincroni: percezione del vicinato $\rightarrow$ computazione locale $\rightarrow$ comunicazione dell'esito $\rightarrow$ comportamento collettivo emergente
    \end{block}
\end{frame}
%
\begin{frame}{Background: Project Emerge}
    \begin{columns}
        \column{0.5\textwidth}
        \centering
        L'obbiettivo del progetto è la realizzazione di un dimostratore open-source per il coordinamento di sciami robotici reali, basato su MacroSwarm + ScaFi     

        \column{0.5\textwidth}
        \begin{block}{Architettura}
            4 moduli -- ognuno responsabile di una determinata mansione -- comunicanti tramite un broker \textbf{MQTT} centralizzato
        \end{block}

        \vspace{1.4em}

        \begin{block}{Vincolo strutturale}
            Ogni componente, software e robot fisici, deve parlare MQTT $\Rightarrow$ nessun supporto per protocolli alternativi
        \end{block}
    \end{columns}
\end{frame}
%
\begin{frame}{Background: connessioni tra i moduli di Emerge}
    \begin{figure}
        \centering
        %\includegraphics[width=0.85\linewidth]{figures/background/connections.pdf}
        \caption{Comunicazione tra i moduli: pubblicazione/sottoscrizione MQTT}
    \end{figure}
\end{frame}
%
\section{Analisi}
%
\begin{frame}{Analisi: obiettivi, vincoli e requisiti}
    \centering
    Obiettivo: rendere il sistema \textbf{protocol-agnostic}, disaccoppiando il protocollo interno (MQTT) da quelli usati dai client esterni

    \vspace{1.5em}

    \begin{columns}
        \column{0.33\textwidth}
        \begin{block}{Vincoli}
            \begin{itemize}
                \item Non modificare i moduli esistenti
                \item MQTT resta il protocollo interno
            \end{itemize}
        \end{block}

        \column{0.33\textwidth}
        \begin{block}{Requisiti funzionali}
            Esporre più protocolli eterogenei, con supporto sia a request/response sia a notifica, ed estendibilità futura
        \end{block}

        \column{0.33\textwidth}
        \begin{block}{Requisiti non funzionali}
            Molte connessioni simultanee, basse latenze, estendibilità senza toccare il nucleo, testabilità
        \end{block}
    \end{columns}
\end{frame}
%
\begin{frame}{Analisi: modello di dominio}
    \begin{figure}
        \centering
        %\includegraphics[width=0.75\linewidth]{figures/analisi/entità.pdf}
        \caption{Schema delle entità di dominio}
    \end{figure}
\end{frame}
%
\section{Progettazione}
%
\begin{frame}{Progettazione: il \textit{domain gateway}}
    \begin{columns}
        \column{0.5\textwidth}
        Un unico modulo aggiuntivo si frappone tra interno ed esterno del sistema, mediando tra broker \textbf{MQTT} e client eterogenei

        \vspace{1em}

        \begin{block}{Vantaggi}
            \begin{itemize}
                \item Moduli esistenti invariati
                \item Nuovi protocolli $\Rightarrow$ si tocca solo il gateway
                \item Supporta sia notifica continua sia request/response (grazie a una \textbf{cache} interna)
            \end{itemize}
        \end{block}

        \column{0.5\textwidth}
        \begin{figure}
            \centering
            %\includegraphics[width=\linewidth]{figures/progettazione/archiettura-sistema.pdf}
            \caption{Architettura del sistema con l'introduzione del gateway}
        \end{figure}
    \end{columns}
\end{frame}
%
\begin{frame}{Progettazione: struttura interna del gateway}
    \begin{figure}
        \centering
        %\includegraphics[width=0.85\linewidth]{figures/progettazione/archiettura-sistema-interno.pdf}
        \caption{Connessioni, bus (Inbound/Outbound) e cache}
    \end{figure}

    \vspace{0.5em}

    \centering
    \textbf{Connessioni} disaccoppiate tramite due \textbf{bus} $\rightarrow$ nuovo protocollo = nuova connessione, nessuna modifica al nucleo
\end{frame}
%
\section{Implementazione}
%
\begin{frame}{Implementazione}
    \begin{columns}
        \column{0.48\textwidth}
        \begin{block}{Stack tecnologico}
            Python 3.13 + \texttt{asyncio}; FastAPI, \texttt{aiocoap}, \texttt{aiomqtt}, Pydantic
        \end{block}

        \vspace{0.8em}

        \begin{block}{Interfacce e classi principali}
            Connection, Bus, Cache, Container, Monitor
        \end{block}

        \column{0.48\textwidth}
        \begin{figure}
            \centering
            %\includegraphics[width=\linewidth]{figures/implementazione/core-overview.pdf}
            \caption{Connection, Bus, Cache orchestrati dal Container}
        \end{figure}

        \vspace{0.8em}

        \centering
        Ogni protocollo rappresenta una sottoclasse indipendente, registrata nel \texttt{Container} 
    \end{columns}
\end{frame}
%
\section{Valutazione}
%
\section{Conclusione}
%
\end{document}